\section{Introduction}\label{sec:introduction}

Global solvers for nonconvex polynomial and mixed-integer nonlinear problems
need lower bounds that are both strong and affordable. Moment-sum-of-squares
(SOS) hierarchies~\cite{lasserre2001-global-optimization-with-polynomials-and,parrilo2003-semidefinite-programming-relaxations-for-semialgebraic}
give a monotone sequence of semidefinite programming (SDP) lower bounds that
converges to the global minimum under standard Archimedean assumptions,
which ensure boundedness through the polynomial generators. The size of a
dense relaxation nevertheless grows quickly with the number of variables.
When the problem has a sparse coupling structure, \emph{correlative sparsity} replaces one dense
moment matrix by one local moment matrix per bag of a junction tree, and
local moments are required to agree on the variables that adjacent bags
share~\cite{waki2006-sums-of-squares-and-semidefinite,lasserre2006-convergent-sdprelaxations-in-polynomial-optimization}.
This keeps the SDP blocks small when the bags are small. Under running
intersection and the corresponding local Archimedean assumptions, the sparse
hierarchy still converges~\cite{lasserre2006-convergent-sdprelaxations-in-polynomial-optimization,grimm2007-structured-sparsity}.

How fast it converges is less well understood. For dense problems on a
box, a series of works based on polynomial kernels established error
$O(r^{-2})$ at relaxation order $r$ for the box preordering (Schmüdgen-type
certificates)~\cite{laurent2023-effective-schmudgen} and, very recently,
$O(\log^3r/r^2)$ for the ordinary box quadratic module (Putinar-type
certificates)~\cite{gribling2026-squared-kernels}. For sparse hierarchies
the published quantitative theorem gives the slower exponent
$2/(w+3)$ for the box preordering, where $w$ is the largest bag
size~\cite{korda2025-convergence-rates-sparsity}. The inverse-square rate
for the sparse preordering was subsequently stated in lecture
slides~\cite{magron2025slides-lorentz,magron2026slides-tenors}; we discuss
this below.

This paper studies the rate of sparse hierarchies, and in particular the
role played by the agreement of separator moments. We consider two
settings. The first is polynomial optimization on a box with a
junction-tree decomposition of the objective. The second, which motivates
much of the paper, attaches to each bag a large vector of private
continuous variables that enter the objective as a convex quadratic and are
constrained by a polytope that may depend affinely on the shared variables.

\subsection{Private convex recourse}\label{sec:intro-recourse}

Many nonconvex models have a small number of coupling decisions and many
continuous decisions that are convex once the coupling decisions are fixed.
As an illustration, consider a network of regions arranged along a tree.
Each region $b$ has a few shared design or price variables $x_{B_b}$, some
of which it shares with neighboring regions, and a dispatch problem in
$p_b$ continuous variables $y_b$ (flows, production levels, storage), with
\[
  \min_{y_b}\ c_b(x)+q_b(x)^{\top}y_b+y_b^{\top}Q_b(x)\,y_b
  \quad\text{subject to}\quad A_by_b\le a_b+E_bx_{B_b},\ \ y_b\in[-1,1]^{p_b}.
\]
The dependence on $x$ can be an arbitrary polynomial (for example, prices
multiplying quantities, or a nonconvex investment cost $c_b$), while for
each fixed $x$ the dispatch problem is a convex quadratic program. If
$E_b\ne0$, the design changes capacities and other right-hand sides; this
is affine recourse.

A standard total-degree SOS hierarchy treats the $v_b+p_b$ variables of a
bag alike, so its moment matrices have $\binom{r+v_b+p_b}{r}$ rows; this is
prohibitive even for moderate $p_b$. Convexity suggests a cheaper
relaxation: let the degree grow only in the shared variables and keep
private moments of degree at most two. The resulting
\emph{private-degree-two hierarchy} (\cref{def:rec-hierarchy}) has local
blocks with $(p_b+1)\binom{r-\abs{I}+v_b}{v_b}$ rows. For two shared
variables per bag, fifty private variables, and $r=6$, the largest block has
$1{,}428$ rows instead of about $4\times10^7$. The question is what accuracy
such a relaxation guarantees and on what that accuracy depends.
Kahl and Henrion use partial lifting of a limited set of nonlinear
variables in geometric reconstruction~\cite{kahl2005-globally-optimal-estimates}.
Guo and Wang keep an SOS-convex decision-variable module at fixed degree
inside a robust polynomial matrix inequality hierarchy whose order grows
in the uncertainty variables~\cite{wang2025-a-moment-sum-of-squares}.
Here we quantify the accuracy of a sparse private-degree-two hierarchy,
with separate rates for fixed domains, affine recourse, and regular
projected multipliers.

\subsection{How separator information can limit convergence}\label{sec:intro-mechanism}

In a sparse hierarchy, each bag communicates with its neighbors through
moments of the shared variables up to degree $2r$. Two actual local
probability measures can therefore use different separator distributions
while agreeing on all the communicated moments. The sharp examples in
\cref{sec:sharpness,sec:recourse} isolate this effect in two bags whose local
minima are $-h(x)$ and $h(x)$ on a one-dimensional separator. In these
paired models the exact local-measure relaxation has gap $2E_{2r}(h)$,
where $E_{2r}(h)$ is the best uniform approximation error by polynomials of
degree at most $2r$~\cite{han2018-local-moment-matching}. The local measures
also give feasible finite-order SDP points, so this identity supplies a
lower bound for their SDP gaps. It need not equal the finite SDP gap or
describe a general sparse model.

For $h(x)=\abs{x}$ the approximation error has exact order $1/r$.
A Lipschitz derivative gives an inverse-square approximation upper bound;
the piecewise quadratic in \cref{thm:quad-sharp} attains that order. A
smoother function, or a polynomial, can be approximated faster. These
examples distinguish the information lost at a separator from errors in
local positivity.

The upper bounds use a common kernel to turn arbitrary feasible truncated
functionals into local laws with exactly matching separator marginals.
For the preordering the kernel is positive under the local cone and
preserves mass. For the ordinary module we first obtain signed exactly
normalized densities, then correct them by one common reference density.
Junction-tree gluing produces a global box law. With private convex
quadratics on fixed polytopes, conditional private means are feasible and
preserve the inverse-square upper bound. Affine recourse makes those means
feasible at a conditional source point rather than the kernel output;
repair at the output costs $O(1/r)$ in general. An active private constraint
can create a corner, but need not do so. Regularity of the projected
multipliers controls the repair more precisely and restores the faster
bounds under the hypotheses of \cref{sec:regularity}.

\subsection{Contributions}\label{sec:intro-contributions}

\Cref{tab:intro-results} summarizes the results. Throughout, $r$ is the
relaxation order, $w$ is the bag width, and the finite-order bounds display
their dependence on the polynomial data, geometry, and regularity parameters.
The box hierarchies use total degree at most $2r$; the recourse hierarchy
uses shared degree at most $2r$ and private degree at most two.

\begin{table}[t]
\centering
\small
\begin{tabular}{@{}>{\raggedright\arraybackslash}p{0.29\textwidth}>{\raggedright\arraybackslash}p{0.19\textwidth}>{\raggedright\arraybackslash}p{0.21\textwidth}>{\raggedright\arraybackslash}p{0.20\textwidth}@{}}
\toprule
Model & Hierarchy & Upper bound & Lower bound / remark\\
\midrule
Box, sparse objective & preordering & $O(r^{-2})$ (\cref{thm:pre}) & $\Omega(r^{-2})$ (\cref{thm:quad-sharp}); rate stated earlier~\cite{magron2025slides-lorentz}\\[3pt]
Box, sparse objective & ordinary module & $\Cf\,O(\log^3r/r^2)$ (\cref{thm:mod}) & $\Omega(r^{-2})$ (\cref{thm:quad-sharp})\\[3pt]
Fixed private polytope, convex QP & private-degree-two & $O(r^{-2})$ (\cref{thm:fixed-recourse}) & needs private quadratic bounds and convexity\\[3pt]
Affine complete recourse & private-degree-two & $O(r^{-1})$ (\cref{thm:affine-recourse}) & $\Omega(r^{-1})$, also total-degree (\cref{thm:affine-sharp})\\[3pt]
Affine recourse, weighted Chebyshev or coordinatewise Lipschitz multipliers & private-degree-two & $O(r^{-2})$ (\cref{thm:reg-cheb}; \cref{thm:reg-holder} with $\beta=1$) & $\Omega(r^{-2})$ (\cref{ex:reg-sharp})\\[3pt]
Polynomial constraints, global Hölder error bound with exponent $\alpha$ & preordering;\newline ordinary module & $O(r^{-\alpha})$;\newline $O((\log r/r)^\alpha)$ (\cref{thm:constr-pre,thm:constr-mod}) & $\Omega(r^{-1})$ possible for $\alpha=1$\\[3pt]
Finite-state labels with box costs & preordering;\newline ordinary module & as for the box (\cref{thm:finite-state,thm:finite-state-mod}) & labels never rounded\\[3pt]
\bottomrule
\end{tabular}
\caption{Summary of the rate results. Upper bounds hold for every feasible
moment family through an explicit rounding construction; lower bounds are
for fixed instances and hold for actual local measures.}
\label{tab:intro-results}
\end{table}

\paragraph{Exactly consistent rounding for the box preordering.}
\Cref{thm:pre} shows that every feasible family of the sparse box
preordering of order $r$ can be rounded to a probability law on the box
whose expected objective exceeds the relaxation objective by at most
$3\Abud/(2m^2+1)$, where $m=\lfloor r/w\rfloor+1$ and $\Abud$ is a weighted
sum of the local Chebyshev coefficients. Hence the gap is at most
$3w^2\Abud/(2r^2)$. We do not claim this rate as new: it was stated for two
overlapping bags in lecture slides of
Magron~\cite{magron2025slides-lorentz,magron2026slides-tenors}, with
attribution to~\cite{korda2025-convergence-rates-sparsity}, whose published
Theorem~6 has the exponent $2/(w+3)$; we have not been able to resolve this
difference. The squared Fejér (Jackson-type) kernel is
classical~\cite{dklerk2017-improved-upper-bounds,laurent2023-effective-schmudgen,korda2025-convergence-rates-sparsity}.
Our proof is self-contained and does not pass through an approximation of
separator value functions. It gives explicit constants, a rounding to a
finite Chebyshev grid (\cref{cor:pre-grid}), and the template for every
later upper bound. A real SOS certificate with the same error follows from
SDP duality (\cref{cor:pre-certificate}).

\paragraph{The ordinary sparse quadratic module.}
The ordinary module uses $v+1$ PSD blocks for a bag of size $v$ instead of
up to $2^v$. The normalized interval-positive kernel used for the
preordering has a tensor certificate involving products of box generators,
which are absent from the ordinary module. Products of globally SOS kernels
are available and provide the starting point for the ordinary construction.
\Cref{thm:mod} transfers the recent dense bound of Gribling, de Klerk and
Vera~\cite{gribling2026-squared-kernels} to the sparse ordinary module: the
gap is at most $\Cf\,O_{w,d}(\log^3r/r^2)$, where $\Cf$ is the sum over bags
of the Chebyshev coefficient norms of the nonconstant parts of the local
objectives and the implicit constant depends only on the width $w$ and the
local degree $d$. After this normalization, the bound does not depend on
the number of bags or on the shape of the tree; the size of the SDP still
grows with the number of bags. The kernel and its SOS normalization are
from~\cite{gribling2026-squared-kernels}. A polynomial kernel that is
nonnegative on the whole source line for every output and has mass
identically one cannot depend on the source variable
(\cref{rem:normalization-obstruction}). Thus this source-dependent SOS kernel
has a mass defect, and its local densities need not have matching separator
marginals. We repair this with an
exactly normalized signed kernel, a pseudomoment Cauchy--Schwarz bound on
the negative parts of the resulting densities, and one common correction
in all bags, which preserves exact consistency. No assumption on an
attained polynomial separator dual is needed, in contrast with the
conditional transfer principle of Gamertsfelder and
Mourrain~\cite{gamertsfelder2025-countable-gmp}. The theorem-specific
contribution is this exactly consistent finite-order transfer, normalized by $\Cf$ and uniform in bag count and tree shape at
fixed width and degree. The slides mention a sparse Putinar rate without
stating its exponent, so we keep the claim separate from priority for a
general sparse rate.

\paragraph{Separator information is the obstruction.}
\Cref{thm:quad-sharp} gives a fixed quadratic in three variables, with two
bags of size two, whose sparse preordering gap is of exact order $r^{-2}$,
while the dense preordering is exact at order two. The lower bound holds even
if every local functional is an actual probability measure, so no
strengthening of local positivity removes it while separators exchange only
moments of degree $2r$. The same witnesses give $\Omega(r^{-2})$ for the
ordinary module. \Cref{prop:exact-orders} computes the first two relaxation
values exactly for both cones; at order two the whole gap is due to
separator information. Qualitative failure of finite sparse exactness with
a nonpolynomial separator is known~\cite{nie2026-sparse-tightness}; the
example adds a sharp rate for a quadratic objective.

\paragraph{Large private convex blocks.}
For a fixed private polytope, \cref{thm:fixed-recourse} proves gap
$O(r^{-2})$ for the private-degree-two hierarchy, with constants that depend
on the shared width and the coefficients but whose dependence on private
dimension enters only through these coefficients. The kernel smooths only
shared variables; private decisions are recovered as conditional means,
which are feasible because the polytope is fixed, and a matrix Jensen
inequality controls the convex cost. Two safeguards are necessary: without
the redundant private quadratic bounds the relaxation can be unbounded
below at every sufficiently high order (\cref{prop:private-bounds-needed}),
and without private convexity the rounding fails even when the objective depends only
on private variables (\cref{prop:convexity-needed}). For this hierarchy
the moment value equals the supremum of certificate values, and every level strictly below it is
certified (\cref{prop:rec-dual}); we do not claim attainment at the
boundary.

\paragraph{Affine recourse and its sharp rate.}
With complete recourse and a fixed private constraint matrix,
\cref{thm:affine-recourse} proves gap $O(r^{-1})$: the conditional private
mean is guaranteed to be feasible at the conditional \emph{source} mean
of the shared variables, but need not be feasible at the output point. A Hoffman error
bound~\cite{hoffman1952-approximate-solutions} repairs it at a cost
proportional to the kernel's transport distance. \Cref{thm:affine-sharp}
shows that $1/r$ is the correct order in general: a two-bag example with one
shared variable $x$, private variables $y_1$ and $y_2$, objective $-xy_1+y_2$,
and private constraints $y_2\ge x$, $y_2\ge-x$ has gap of exact order
$1/r$, for the private-degree-two hierarchy and, separately, for the
total-degree sparse preordering. Merging the two bags, or splitting at the
active-set change, gives an exact full-preordering certificate at low order.

\paragraph{Regular multipliers restore the inverse-square rate.}
The sharp affine example has a corner in its separator value function and
a discontinuity in its projected multipliers.
\Cref{thm:reg-cheb,thm:reg-holder} show that the rate $r^{-2}$ returns, for
the same unchanged hierarchy, if the \emph{projected} optimal multipliers---the
part of the dual solution that acts on the shared variables---have
absolutely summable
weighted tensor Chebyshev coefficients, or each component is a Lipschitz
function of its corresponding shared coordinate, the case $\beta=1$ of
the Hölder theorem. Here the $i$th component has the form
$\ell_{b,i}(x_i)$, with $\ell_{b,i}$ univariate.
More generally, coordinatewise Hölder dependence of order
$0<\beta\le1$ gives $O(r^{-(1+\beta)})$. A key step bounds pseudomoments
of Chebyshev polynomials whose total degree exceeds $r$ (\cref{lem:half-degree}). Strict
convexity of the private problem does not imply the required regularity
(\cref{ex:strict-convex-insufficient}); classical parametric quadratic
programming gives sufficient
conditions~\cite{tondel2003-an-algorithm-for-multi-parametric,baotic2016-gradient-value-function}. We do not
cover arbitrary multivariate Lipschitz projected multipliers.

\paragraph{Extensions.}
\Cref{sec:extensions} adds polynomial constraints under a \emph{global}
Hölder error bound with exponent $\alpha\in(0,1]$. The rounded law has small
expected squared constraint violation and is then repaired onto the feasible
set, giving $O(r^{-\alpha})$ for box-preordering products times single
constraints and $O((\log r/r)^\alpha)$ for the ordinary constrained module
(\cref{thm:constr-pre,thm:constr-mod}). The ordinary-module proof bounds
expected squared violation at scale $O(\log^2r/r^2)$ through a
polynomial displacement certificate, then uses global geometric repair.
The same rates hold for real certificates with arbitrarily small positive
slack, without strict feasibility (\cref{prop:constr-dual}); boundary
attainment is not established. A global error bound can be much worse than
separate local bounds. Finite-state variables with bag-local label restrictions are handled without rounding
any label, for both cones (\cref{thm:finite-state,thm:finite-state-mod}).
Finally, \cref{sec:certificates} turns the box bounds into exact rational
certificates with any rational positive slack, with encoding length
polynomial in the expanded SDP size and rational input length, including
the slack. Positive slack beyond the
proved finite-order error supplies the real cone membership needed for this
specialization of established rational SOS
methods~\cite{peyrl2008-computing-sum-of-squares-decompositions,davis2024-rational-dual-certificates-for-weighted}.

\subsection{Discretization and computational interpretation}\label{sec:intro-limits}

The rounding theorems bound the increase in objective when a feasible moment
family is converted into a feasible law, and hence bound the error of the
specified relaxation. In the private convex setting this comparison is
one-sided, because taking conditional means can improve the objective.
The accuracy bounds alone do not give numerical running times: SDP size
grows with bag count, and geometric and coefficient constants can be large.
Discretization also exploits bounded interaction width: Bienstock and
Mu\~noz develop approximate linear programming formulations for polynomial
optimization with bounded intersection treewidth~\cite{bienstock2018-lp-formulations-for-polynomial-optimization}.
For our box problems, evaluating the objective on a Chebyshev grid and
minimizing by dynamic programming over the junction tree gives error
$O(r^{-2})$ with $O(r)$ nodes per coordinate (\cref{cor:pre-grid}).
With exact local
quadratic programs, the analogous grid bound is $O(r^{-2})$ for fixed
private domains and $O(1/r)$ for general affine complete recourse
(\cref{rem:grid-baseline}). If a proved grid error is $\varepsilon_r$ and
the grid optimum is $G_r$, then $G_r-\varepsilon_r\le\fstar\le G_r$ also
provides a lower bound. The hierarchy results instead control the specified
sparse SDP, every feasible moment point, and the associated algebraic
certificate cone. Their accuracy bounds do not establish a generic runtime
or practical speedup.

\subsection{Organization}\label{sec:intro-organization}

\Cref{sec:setting} fixes notation, the three hierarchies, and standard tools.
\Cref{sec:kernels} proves the preordering rounding theorem and introduces
the construction used throughout. \Cref{sec:ordinary} treats the ordinary
module, and \cref{sec:sharpness} gives the matching lower bound and the
exact small-order values. \Cref{sec:recourse} studies fixed and affine
private recourse, and \cref{sec:regularity} the regular-multiplier rates.
\Cref{sec:extensions} covers polynomial constraints and finite-state
variables, and \cref{sec:certificates} rational certificates.
\Cref{sec:discussion} discusses limitations and open problems. The appendix
proves the failure result when private quadratic bounds are omitted.
